\documentclass[11pt,a4paper]{article}

\usepackage[utf8]{inputenc}
\usepackage[T1]{fontenc}
\usepackage[spanish,es-nodecimaldot]{babel}
\usepackage[a4paper,margin=2.2cm]{geometry}
\usepackage{xcolor}
\usepackage{listings}
\usepackage{hyperref}
\usepackage{enumitem}
\usepackage{fancyhdr}

\definecolor{ink}{HTML}{1F2937}
\definecolor{accent}{HTML}{176B87}
\definecolor{codebg}{HTML}{F3F6F8}
\definecolor{codeframe}{HTML}{D5E0E5}

\hypersetup{
  colorlinks=true,
  linkcolor=accent,
  urlcolor=accent,
  pdftitle={Periodicity Workbench: guía de inicio local},
  pdfauthor={}
}

\lstdefinestyle{terminal}{
  backgroundcolor=\color{codebg},
  frame=single,
  rulecolor=\color{codeframe},
  basicstyle=\ttfamily\small,
  columns=fullflexible,
  keepspaces=true,
  breaklines=true,
  showstringspaces=false,
  framexleftmargin=5pt,
  framexrightmargin=5pt,
  aboveskip=7pt,
  belowskip=9pt
}
\lstset{style=terminal}

\pagestyle{fancy}
\fancyhf{}
\lhead{Periodicity Workbench}
\rhead{Guía de inicio local}
\cfoot{\thepage}
\setlength{\headheight}{14pt}

\title{\textcolor{accent}{Periodicity Workbench}\\[3pt]
\large Guía breve para instalar y ejecutar en local}
\author{}
\date{}

\begin{document}
\maketitle
\vspace{-1.2em}

Esta guía explica cómo instalar Git y Python, descargar el proyecto e iniciar la aplicación en Windows o Linux. Streamlit y el resto de las bibliotecas se instalan desde el archivo \texttt{requirements.txt} del proyecto.

\section{Instalar Git y Python}

\subsection{Windows}
\begin{enumerate}[leftmargin=*]
  \item Instala \href{https://git-scm.com/install/win}{Git for Windows}.
  \item Instala Python desde \href{https://www.python.org/downloads/}{python.org}. Python 3.12 o posterior es una buena opción.
  \item Abre PowerShell y comprueba que ambos estén disponibles:
\end{enumerate}

\begin{lstlisting}
git --version
py --version
\end{lstlisting}

\subsection{Linux}
Abre una terminal e instala Git, Python, \texttt{pip} y el soporte para entornos virtuales. Elige los comandos de tu distribución.

\textbf{Ubuntu o Debian:}
\begin{lstlisting}
sudo apt update
sudo apt install git python3 python3-venv python3-pip
\end{lstlisting}

\textbf{Fedora:}
\begin{lstlisting}
sudo dnf install git python3 python3-pip
\end{lstlisting}

Comprueba la instalación:
\begin{lstlisting}
git --version
python3 --version
\end{lstlisting}

\section{Descargar el proyecto}

En PowerShell o en la terminal de Linux, ve a la carpeta donde quieras guardar el proyecto y ejecuta:

\begin{lstlisting}
git clone https://github.com/pereblay/periodicity-workbench.git
cd periodicity-workbench
\end{lstlisting}

\section{Crear un entorno e instalar las dependencias}

Un entorno virtual mantiene las bibliotecas del proyecto separadas de las del resto del equipo.

\textbf{Windows PowerShell:}
\begin{lstlisting}
py -m venv .venv
.\.venv\Scripts\Activate.ps1
python -m pip install --upgrade pip
pip install -r requirements.txt
\end{lstlisting}

\textbf{Linux:}
\begin{lstlisting}
python3 -m venv .venv
source .venv/bin/activate
python -m pip install --upgrade pip
pip install -r requirements.txt
\end{lstlisting}

La instalación puede tardar unos minutos. El archivo instala Streamlit y todas las bibliotecas necesarias para la aplicación.

\section{Arrancar la aplicación}

Con el entorno virtual activado y dentro de la carpeta \texttt{periodicity-workbench}, ejecuta:

\begin{lstlisting}
python -m streamlit run streamlit_app.py
\end{lstlisting}

Se abrirá el navegador. Si no se abre automáticamente, visita \href{http://localhost:8501}{http://localhost:8501}. Para detener la aplicación, vuelve a la terminal y pulsa \textbf{Ctrl+C}.

\section{La próxima vez}

No hace falta volver a clonar el proyecto ni instalar las bibliotecas. Entra en la carpeta del proyecto, activa el entorno y arranca la aplicación.

\textbf{Windows PowerShell:}
\begin{lstlisting}
cd periodicity-workbench
.\.venv\Scripts\Activate.ps1
python -m streamlit run streamlit_app.py
\end{lstlisting}

\textbf{Linux:}
\begin{lstlisting}[language=bash]
cd periodicity-workbench
source .venv/bin/activate
python -m streamlit run streamlit_app.py
\end{lstlisting}

\section*{Si Linux no encuentra \texttt{venv}}

En Ubuntu o Debian, instala el paquete que falta y vuelve a crear el entorno:
\begin{lstlisting}[language=bash]
sudo apt install python3-venv
\end{lstlisting}

\vfill
\noindent\small Documentación oficial: \href{https://docs.streamlit.io/get-started/installation/command-line}{instalación de Streamlit}, \href{https://git-scm.com/install/win}{Git para Windows} y \href{https://git-scm.com/install/linux}{Git para Linux}.

\end{document}
